%=============================================================================!
% 6.2  THE ACTION AND HAMILTON'S PRINCIPLE                                    !
%=============================================================================!
\section{The action and Hamilton's principle}
\label{sec:la-action}

Newton's laws work forwards from a starting state, one instant to the
next. There is a second way to single out the true motion: compare whole
paths that start at one place at one time and end at another place at a
later time, and pick the one that makes a certain number stationary. This
section defines that number and tests the statement on a thrown ball;
\cref{sec:la-euler} shows that it is equivalent to Newton's laws.

\subsection*{The Lagrangian and the action}

The Lagrangian of a system is its kinetic energy minus its potential
energy,
\begin{equation}
   \Lag(q, \dot{q}) = K - U ,
   \label{eq:la-lagrangian}
\end{equation}
written, as in \cref{sec:la-coordinates}, in terms of the generalised
coordinates and velocities, $q$ standing for all of $q_1, \dots, q_n$ and
$\dot{q}$ for all of $\dot{q}_1, \dots, \dot{q}_n$. For a path $q(t)$ from the time $t_1$ to the
time $t_2$, its action is the integral
\begin{equation}
   \Act = \int_{t_1}^{t_2}\Lag\bigl(q(t), \dot{q}(t)\bigr)\,\dd t ,
   \label{eq:la-action}
\end{equation}
one number for the whole path, in joule seconds. Each path, true or not,
has an action: put its $q(t)$ and $\dot{q}(t)$ into $\Lag$ and integrate.

\emph{Hamilton's principle} states that, among all paths that start at
$q_{\mathrm{A}}$ at the time $t_1$ and end at $q_{\mathrm{B}}$ at the time
$t_2$, the path the system actually follows makes the action stationary:
changing the path from $q(t)$ to $q(t) + \zeta\eta(t)$, with $\eta$ any bump
that is zero at $t_1$ and $t_2$ and $\zeta$ a small number, changes $\Act$
by no amount proportional to $\zeta$ itself, only by amounts of order
$\zeta^2$ (\cref{sec:mo-taylor}), as at the bottom of a valley or the top
of a hill of $\Act$ against $\zeta$.

\subsection*{A thrown ball}

A ball of \SI{0.145}{\kilogram} is thrown straight up at
\SI{12}{\metre\per\second}, as in Chapters~1 to~3. Its height is $x(t) =
12t - \frac12 gt^2$ (\cref{sec:mo-integrating}), and after
\SI{2}{\second} it is \SI{4.387}{\metre} up and falling. Now invent other
paths between the same two events (an event is a place at a given time), starting at $x = 0$ at $t = 0$ and
ending at \SI{4.387}{\metre} at \SI{2}{\second}: add to the true height
$\zeta\eta(t)$, with $\eta(t) = \sin(\pi t/\tau)$ and $\tau =
\SI{2}{\second}$, a bump that is zero at both ends. Each trial path has its
own action, with $\Lag = \frac12 m\dot{x}^2 - mgx$, and the function
\texttt{action} of \texttt{mdlab} computes it by splitting the path into
short straight steps (\cref{fig:la-action}). For the true path, $\zeta =
0$, the integral can also be done by hand. With $\dot{x} = 12 - gt$,
\begin{align*}
   \Lag &= \tfrac12 m\,(12 - gt)^2 - mg\,\bigl(12t - \tfrac12 gt^2\bigr)\\
   &= 72m - 12mg\,t + \tfrac12 mg^2t^2 - 12mg\,t + \tfrac12 mg^2t^2\\
   &= 72m - 24mg\,t + mg^2t^2 ,
\end{align*}
expanding the square and the bracket, then collecting like terms,
and, integrating each term by the power rule and the fundamental theorem
(\cref{sec:mo-integrating}),
\begin{equation*}
   \Act = \Bigl[72m\,t - 12mg\,t^2 + \tfrac13 mg^2t^3\Bigr]_0^{2}
   = 20.8800 - 68.2543 + 37.1859 = \SI{-10.1884}{\joule\second} ,
\end{equation*}
as \texttt{action} also finds. The paths with $\zeta = \pm\SI{1}{\metre}$
have \SI{-10.0095}{\joule\second}, and those with $\pm\SI{2}{\metre}$ have
\SI{-9.4729}{\joule\second}. The excess grows as $\zeta^2$: \num{0.1789}
and \num{0.7155}, four times as much for twice the change, even for bumps
of metres, which are not small; the next paragraph shows why. Among these
paths the true one has the least action.

\begin{figure}[htb]
\centering
\includegraphics{ch06_lagrange/action}
\caption{Trial paths of a ball of \SI{0.145}{\kilogram} thrown up at
\SI{12}{\metre\per\second}, all starting at $x = 0$ at $t = 0$ and ending
\SI{4.387}{\metre} up at \SI{2}{\second}. (a) The true path (teal) and the
trial paths with $\zeta = \pm1$ (grey) and $\pm\SI{2}{\metre}$ (ochre).
(b) The action of each trial path against $\zeta$, least at $\zeta = 0$.}
\label{fig:la-action}
\end{figure}

The excess can be found exactly. The change of path changes the velocity
by $\zeta\dot\eta$ and the height by $\zeta\eta$, so
\begin{align*}
   \Act(\zeta) &= \int_0^\tau\Bigl[\tfrac12 m(\dot{x} + \zeta
      \dot\eta)^2 - mg\,(x + \zeta\eta)\Bigr]\dd t\\
   &= \int_0^\tau\Bigl[\bigl(\tfrac12 m\dot{x}^2 - mgx\bigr)
      + \zeta\bigl(m\dot{x}\dot\eta - mg\eta\bigr)
      + \zeta^2\,\tfrac12 m\dot\eta^2\Bigr]\dd t\\
   &= \Act(0) + \zeta\int_0^\tau\bigl(m\dot{x}\dot\eta - mg\eta\bigr)\,
      \dd t + \zeta^2\,\tfrac12 m\int_0^\tau\dot\eta^2\,\dd t ,
\end{align*}
expanding the square, grouping by powers of $\zeta$, and splitting the
integral.
There are no higher powers of $\zeta$, which is why even large bumps
follow $\zeta^2$ exactly. \Cref{sec:la-euler} shows that the middle
integral vanishes for the true path and any bump $\eta$. The last needs
$\dot\eta$ and the identity $\cos^2\theta = \frac12 + \frac12\cos2\theta$
of \cref{sec:os-energy}:
\begin{align*}
   \int_0^\tau\dot\eta^2\,\dd t
   &= \Bigl(\frac{\pi}{\tau}\Bigr)^2\int_0^\tau\cos^2\frac{\pi t}{\tau}\,
      \dd t
      && \text{$\dot\eta = \frac{\pi}{\tau}\cos\frac{\pi t}{\tau}$ by the
         chain rule,}\\
   &= \Bigl(\frac{\pi}{\tau}\Bigr)^2\int_0^\tau\Bigl(\tfrac12
      + \tfrac12\cos\frac{2\pi t}{\tau}\Bigr)\dd t
      && \text{the identity, with $\theta = \pi t/\tau$,}\\
   &= \Bigl(\frac{\pi}{\tau}\Bigr)^2\Bigl[\frac{t}{2}
      + \frac{\tau}{4\pi}\sin\frac{2\pi t}{\tau}\Bigr]_0^\tau
      && \text{the fundamental theorem,}\\
   &= \Bigl(\frac{\pi}{\tau}\Bigr)^2\frac{\tau}{2} = \frac{\pi^2}{2\tau}
      && \text{as $\sin 2\pi = \sin 0 = 0$.}
\end{align*}
The term in $\zeta^2$ is therefore $\zeta^2\times\frac12 m\times\pi^2/(2
\tau) = \zeta^2\,m\pi^2/(4\tau)$, and $m\pi^2/(4\tau) = 0.145\times9.8696/8
= \SI{0.17889}{\joule\second\per\metre\squared}$, the excess found above
for $\zeta = \SI{1}{\metre}$.

\subsection*{When the action is not least}

The true path is not always the one of least action. Take the block of
\cref{sec:os-circle}, with $m = \SI{0.5}{\kilogram}$ and $k =
\SI{50}{\newton\per\metre}$, so that $\omega = \sqrt{k/m} =
\SI{10}{\radian\per\second}$, oscillating for $\tau = \SI{0.5}{\second}$,
longer than half its period, $\pi/\omega = \SI{0.3142}{\second}$. Now
$\Lag = \frac12 m\dot{x}^2 - \frac12 kx^2$, and the potential energy of the
changed path, $\frac12 k(x + \zeta\eta)^2 = \frac12 kx^2 + \zeta kx\eta +
\zeta^2\,\frac12 k\eta^2$, brings a second term in $\zeta^2$. The term in
$\zeta$ again vanishes for the true path (\cref{sec:la-euler}), whatever
its amplitude, so for $\eta = \sin(\pi t/\tau)$ the excess is
\begin{align*}
   \Act(\zeta) - \Act(0)
   &= \zeta^2\Bigl(\tfrac12 m\int_0^\tau\dot\eta^2\,\dd t
      - \tfrac12 k\int_0^\tau\eta^2\,\dd t\Bigr)
      && \text{expanding as for the ball,}\\
   &= \zeta^2\Bigl(\tfrac12 m\Bigl(\frac{\pi}{\tau}\Bigr)^2\frac{\tau}{2}
      - \tfrac12 k\,\frac{\tau}{2}\Bigr)
      && \text{as above, with $\sin^2\theta = \tfrac12 - \tfrac12\cos2\theta$,}\\
   &= \zeta^2\,\frac{\tau}{4}\Bigl(m\Bigl(\frac{\pi}{\tau}\Bigr)^2 - k\Bigr)
      && \text{taking out $\tau/4$.}
\end{align*}
The bracket is negative whenever $\pi/\tau < \sqrt{k/m} = \omega$, that is
whenever $\tau$ is longer than $\pi/\omega$: here it is $0.125\times(0.5
\times39.48 - 50)\,\zeta^2 = -3.783\,\zeta^2$, and the action falls. For the
faster bump $\sin(2\pi t/\tau)$ the same steps with $2\pi$ in place of $\pi$
give $0.125\times(0.5\times157.91 - 50)\,\zeta^2 = +3.620\,\zeta^2$, and the
action rises. Against the sizes of the two bumps the true path sits at a
saddle of the action, like the saddle points of \cref{sec:en-gradient}:
stationary, which is all Hamilton's principle claims.

\begin{keyidea}
The Lagrangian is $\Lag = K - U$ and the action of a path is $\Act =
\int\Lag\,\dd t$. Among all paths with the same ends in space and time,
the true motion makes the action stationary: a small change of path, of
size $\zeta$, changes it only by an amount of order $\zeta^2$.
\end{keyidea}

\notebook{06}{bend the path of the ball and watch its action}

\begin{selfcheck}
For a body with no forces on it, $U = 0$, the true path between two
events is a straight line at constant speed. Using the expansion above,
why is its action less than that of every other path with the same ends?
\end{selfcheck}
